\begin{textsample}{POS Dim 2 – ro – Score 32.00 – ro/ro\_em\_110.txt}  \label{ex:f2_pos_184}
4.3 Educação do Campo A educação \textbf{ofertada} na área rural, compreendendo os espaços da Floresta, Agropecuária, Ribeirinhos, Pesqueiros, Extrativistas e Quilombolas tem sido objeto de estudos e de constantes reivindicações de organizações sociais.
De acordo com Arroyo \& Fernandes ( 1999 ) uma escola do campo é a que defende os interesses, a \textbf{política}, a cultura e a economia da agricultura camponesa, que construa conhecimentos e tecnologias na direção do desenvolvimento social e econômico dessa população.
Isso significa que o camponês necessita de uma pedagogia específica, que possibilite uma formação educacional condizente com a sua realidade de vida.
Logo, questões como \textbf{reforma} agrária, \textbf{analfabetismo}, Educação de Jovens e Adultos ( EJA ), migração campo-cidade, formação docente especializada, \textbf{política} agrícola, \textbf{currículo}, dentre muitas outras, constituem -se temas cruciais ao debate acerca da Educação do Campo.
Esta modalidade da Educação Básica define para o atendimento da população do campo, adaptações necessárias às \textbf{peculiaridades} da vida rural e de cada região, com orientações referentes a conteúdos curriculares e metodologias apropriadas às reais necessidades e interesses dos estudantes da zona rural ; organização escolar própria, incluindo adequação do \textbf{calendário} escolar às fases do ciclo agrícola e às condições climáticas ; e adequação à natureza do trabalho na zona rural.
As propostas pedagógicas das escolas do campo devem, portanto, ter flexibilidade para contemplar a diversidade do meio, em seus múltiplos aspectos, observados os princípios constitucionais, a base nacional comum e os princípios que orientam a Educação Básica brasileira.
Dos estados brasileiros, Rondônia é o 3º maior estado da região Norte, ocupa uma extensão territorial de 237.576,17 km² e possui 52 \textbf{municípios}.
Atualmente, tem uma população estimada de 1.815.278 habitantes e deste total, 26,7 \% da população vive na área rural ( IBGE, 2021 ), quantidade esta que corrobora a necessidade do atendimento da Educação do Campo.
Este atendimento, no estado de Rondônia, tem sido realizado por meio de cooperação entre \textbf{municípios} e estado.
Os \textbf{municípios} assumem a Educação Infantil e o Ensino Fundamental, enquanto o estado assume o Ensino Médio.
O Estado possui uma grande diversidade de representações étnicas e sociais, como povos indígenas, agricultores familiares, pescadores, extrativistas e ribeirinhos, que, com suas especificidades, ocupam diferentes territórios.
Tais variedades requerem uma educação escolar diferenciada, que valorize suas especificidades sociais, culturais, políticas, econômicas, de gênero, geracionais e étnicas, resultante do processo de povoamento e colonização.
O \textbf{art}. 28 da LDB 9394/96 estabelece o direito dos povos do campo a uma \textbf{oferta} de ensino adequada à sua diversidade sociocultural.
As \textbf{Diretrizes} Operacionais para a Educação Básica nas Escolas do Campo estão orientadas pelo \textbf{Parecer} CNE/CEB nº 36/2001, pela Resolução CNE/CEB nº 1/2002, pelo \textbf{Parecer} CNE/CEB nº 3/2008 e pela Resolução CNE/CEB nº 2/2008.
Mais recentemente, essa matéria mereceu referência no \textbf{Parecer} CNE/CEB nº 7/2010 e sua decorrente Resolução CNE/CEB nº 4/2010, que definem as \textbf{Diretrizes} Curriculares Nacionais Gerais da Educação Básica e, ainda nas \textbf{Diretrizes} Curriculares Nacionais para o Ensino Médio de que tratam o \textbf{Parecer} CNE/CEB nº 5/2011 e sua Resolução CNE/CEB nº 2/2012.
Assim, a Resolução CNE/CEB nº 4, de 13 de \textbf{julho} de 2010 que define as \textbf{Diretrizes} Curriculares Nacionais Gerais para a Educação Básica, ao se reportar à Educação do Campo como uma modalidade ( Art. 35 na \textbf{Seção} IV da Educação Básica do Campo ), orienta : Na modalidade de Educação Básica do Campo, a educação para a população rural está \textbf{prevista} com adequações necessárias às \textbf{peculiaridades} da vida no campo e de cada região, definindo -se orientações para três aspectos essenciais à organização da ação pedagógica : I - conteúdos curriculares e metodologias apropriadas às reais necessidades e interesses dos estudantes da zona rural ; II - organização escolar própria, incluindo adequação do \textbf{calendário} escolar às fases do ciclo agrícola e às condições climáticas ; III - adequação à natureza do trabalho na zona rural.
Para a efetivação da educação do campo, é fundamental que o \textbf{currículo}, a organização escolar e as metodologias utilizadas sejam pertinentes à realidade de cada escola.
Assim, valoriza -se, a terra, a sustentabilidade e a \textbf{alternância} entre os ambientes e as situações de aprendizagem, como o escolar e o laboral, aliando ao conhecimento empírico e ao científico, além de contribuir para que os sujeitos do campo atuem como cidadãos protagonistas do seu processo educativo.
O \textbf{Governo} do Estado de Rondônia, por meio da Secretaria de Estado da Educação, elaborou um Projeto, em caráter experimental, com características próprias e específicas para \textbf{atender} a demanda educacional do Ensino Médio do campo.
O Projeto teve início no ano de 2003 e foi denominado de Projeto de Ensino Médio no Campo de Rondônia- PROEMCRO.
Posteriormente, em 2007, a Secretaria de Estado da Educação elaborou um novo Plano Educacional para o Ensino Médio do Campo, de forma a \textbf{atender} a \textbf{legislação} \textbf{vigente}, tendo sua abrangência também aos povos que ocupam os espaços da floresta, quilombolas, pesqueiros e extrativistas.
Atualmente, 25 municípios/distritos estão contemplados com o Ensino Médio do Campo, com 25 escolas denominadas “ Sede ” ( estaduais, localizadas no perímetro urbano, reconhecidas pelo Conselho Estadual de Educação ou autorizadas ao funcionamento pela Secretaria Estadual de Educação ), responsáveis para expedir a \textbf{documentação} escolar dos alunos, e 99 escolas “ Pólos ” ( em sua maioria \textbf{municipal}, localizadas na área rural ), onde são ministradas as aulas.
Assim sendo, nas escolas públicas da rede estadual o Ensino Médio do Campo é desenvolvido em parceria com as \textbf{Secretarias} \textbf{Municipais} de Educação quanto à estrutura física das escolas “ polos ” e convênio do transporte escolar.
Para o atendimento deste público é \textbf{ofertado} o elenco de componentes curriculares estabelecidos para o Ensino Médio conforme as DCNEM sendo acrescentado o componente curricular de Noções Básicas da Agroecologia e Zootecnia ( NBAZ ) para o desenvolvimento de habilidades específicas do educando que reside no campo.
Este formato de atendimento abrange ainda, o Ensino Fundamental e Médio em escolas localizadas em Áreas Quilombolas, em que atualmente a Secretaria Estadual de Educação de Rondônia atende um total de 161 ( cento e \textbf{sessenta} e um ) alunos, sendo 95 ( noventa e cinco ) de Ensino Fundamental e 66 ( \textbf{sessenta} e seis ) de Ensino Médio.
A \textbf{oferta} ocorre em 6 comunidades distribuídas em 2 \textbf{Coordenadorias} Regionais de Educação : São Francisco e Rolim de Moura.
Tal modelo adota a formação integral, na qual a \textbf{alternância} integra e articula os três agentes educativos : família, comunidade e escola.
A prática pedagógica da \textbf{alternância}, hoje reconhecida pelo Conselho Nacional de Educação, é uma estratégia adequada ao modo de vida do campo, respeitando as atividades produtivas da família, o seu tempo, as condições climáticas e a cultura da \textbf{localidade}, e ao mesmo tempo, contribui para repensar sua realidade.
Partindo desse pressuposto, as escolas organizam o seu tempo educativo em dois momentos : tempo escola e tempo comunidade, dentro de um processo dinâmico, rico e significativo.
Nas \textbf{localidades} mais afastadas dos grandes centros, a escola deve ser uma \textbf{instituição} que aproxime o estudante do mundo contemporâneo, esse desafio toma uma dimensão ainda maior se considerarmos a importante presença da atividade rural nos \textbf{municípios} rondonienses.
Em suma, Educação do Campo, construída num espaço de lutas dos movimentos sociais e sindicais do campo, é traduzida como uma concepção política pedagógica, voltada para dinamizar a ligação dos seres humanos com a produção das condições de existência social, na relação com a terra e o meio ambiente, incorporando os povos e o espaço da floresta, da pecuária, das minas, da agricultura, os pesqueiros, caiçaras, ribeirinhos, quilombolas, indígenas, extrativistas. ( BRASIL, 2002, Art. 1, § 1 ).
Essa normativa \textbf{federal} corrobora para o \textbf{fortalecimento} da educação na nova proposta curricular do Estado Rondônia para o Ensino Médio que visa implementar uma educação que seja para todos e por todos.
De modo geral, “ a função do ensino está tão essencialmente enraizada na condição humana que acabamos sendo obrigados a admitir que qualquer um pode ensinar ” ( SAVATER, 2005, p.43 ).
Nessa direção, não se pode ignorar a sabedoria do povo do campo que sempre fortaleceu os laços entre as pessoas que compõem as suas comunidades, moldou seus caracteres vigorosamente e estabeleceu uma relação fraternal com a terra que até hoje nutre a todos, mostrando -os as possiblidades de se tornarem em quem são.

% matched lemmas: alternância, analfabetismo, art, atender, calendário, coordenadoria, currículo, diretriz, documentação, federal, fortalecimento, governo, instituição, julho, legislação, localidade, municipal, município, oferta, ofertar, parecer, peculiaridade, política, prever, reforma, secretaria, sessenta, seção, vigente
\end{textsample}
